\documentclass[11pt]{article}

% 基础包
\usepackage[utf8]{inputenc}
\usepackage[T1]{fontenc}
\usepackage{amsmath,amssymb,amsfonts}
\usepackage{booktabs}
\usepackage{graphicx}
\usepackage{hyperref}
\usepackage{geometry}
\geometry{margin=1in}
\usepackage{xcolor}

% 标题设置
\title{\textbf{Attribution-Guided Online Distillation (AGOD):}\\
\large Training on the Shift, Not Just the Data}

\author{Daily Iteration Prototype}
\date{\today}

\begin{document}
\maketitle

\begin{abstract}
Real-world multi-modal systems operate under continuous distribution shift. Traditional distillation pipelines treat all modalities uniformly, wasting compute on stable modalities while under-correcting drifted ones. In this work, we propose \textbf{Attribution-Guided Online Distillation (AGOD)}, a framework that reframes modality-specific gap decomposition as a \emph{clever covariate} to dynamically control the distillation objective. At each online step, we quantify the Modality-Specific Gap (MSG) via Random-Forest domain classifiers and meta-learner pseudo-outcome risk. The resulting state vector is mapped to dynamic modality weights, effectively routing distillation capacity toward the drifting modality. We instantiate AGOD on ChronoBerg and demonstrate that attribution-guided routing outperforms static and covariate-only distillation. This work positions a new branch of online learning where the update signal is a distribution-shift statistic.
\end{abstract}

\section{Introduction}
\label{sec:intro}

Multi-modal representation learning (CLIP, CLAP) has become the standard backbone for retrieval, recommendation, and generation. In production, these models are distilled into lightweight students for latency and cost. However, production data is inherently non-stationary: user mixtures shift, audio environments change, and text semantics evolve.

Standard distillation assumes stationarity. When data drifts, practitioners either (i) retrain everything (expensive, risks catastrophic forgetting) or (ii) ignore the drift (silent degradation). Both are unacceptable in an online setting.

We argue that \textbf{the distillation objective itself should be a function of the distribution shift}. Rather than treating attribution as a post-hoc diagnostic, we promote it to a \emph{control signal}: a clever covariate that modulates which modality, which sample, and which token receives distillation capacity at each online step.

\paragraph{Core Philosophy.}
We do not train on the data; we train on the \emph{shift} of the data. Modality imbalance is no longer a static hyperparameter; it is a dynamic control signal.

\section{Methodology: AGOD}
\label{sec:method}

\subsection{Modality-Specific Gap (MSG) Decomposition}

Let \(\mathcal{M} = \{\text{audio}, \text{image}, \text{text}\}\). At time step \(t\), we observe a new batch \(\mathcal{D}_t\) and maintain a reference batch \(\mathcal{D}_0\). For each modality \(m \in \mathcal{M}\), we compute:
\begin{align}
    \text{AUC}_m^{(t)} &= \text{RF-Domain}(\mathbf{X}_m^{\mathcal{D}_0}, \mathbf{X}_m^{\mathcal{D}_t}), \\
    \text{VIMP}_m^{(t)} &= \text{OOB-VIMP}(\text{RF-Domain}), \\
    \text{PO-risk}_m^{(t)} &= \text{CFPerm}(\mathbf{X}_m^{\mathcal{D}_0}, Y^{\mathcal{D}_0}, \mathbf{X}_m^{\mathcal{D}_t}, Y^{\mathcal{D}_t}).
\end{align}

We define the \textbf{Modality-Specific Gap} as:
\begin{equation}
    g_m^{(t)} = \text{Normalize}\left( \text{AUC}_m^{(t)} \cdot \text{VIMP}_m^{(t)} + \gamma \cdot \text{PO-risk}_m^{(t)} \right),
\end{equation}
and the state vector \(\mathbf{g}^{(t)} = [g_{\text{audio}}^{(t)}, g_{\text{image}}^{(t)}, g_{\text{text}}^{(t)}]\).
\textit{Remark:} \(\mathbf{g}^{(t)}\) is a distance-based statistic quantifying how far each modality has moved, not a causal claim.

\subsection{Clever Covariate Construction}

We treat \(\mathbf{g}^{(t)}\) as a clever covariate. The dynamic modality weights are:
\begin{equation}
    \boldsymbol{\alpha}^{(t)} = \text{Softmax}\left( \frac{\mathbf{g}^{(t)}}{\tau} \right),
\end{equation}
where \(\tau\) is a temperature controlling routing sharpness. When \(\text{AUC}_{\text{audio}}^{(t)}\) spikes, \(\alpha_{\text{audio}}^{(t)} \to 1\), and distillation capacity is routed to audio.

\subsection{Attribution-Guided Distillation Objective}

Let \(f_T\) be the frozen teacher (e.g., CLAP) and \(f_S\) the student. The online distillation loss at step \(t\) is:
\begin{equation}
    \mathcal{L}^{(t)} = \mathcal{L}_{\text{global}}^{(t)} + \sum_{m \in \mathcal{M}} \alpha_m^{(t)} \mathcal{L}_m^{(t)} + \lambda \Omega(\boldsymbol{\alpha}^{(t)}),
\end{equation}
where:
\begin{itemize}
    \item \(\mathcal{L}_{\text{global}}^{(t)}\): global similarity-matrix matching (KL divergence).
    \item \(\mathcal{L}_m^{(t)}\): modality-specific local alignment (slice-token or patch-token similarity).
    \item \(\Omega(\boldsymbol{\alpha}^{(t)})\): temporal smoothness regularizer to prevent weight oscillation.
\end{itemize}

\subsection{Online Adaptation Loop}
\begin{enumerate}
    \item \textbf{Monitor}: receive batch \(\mathcal{D}_t\), compute \(\mathbf{g}^{(t)}\).
    \item \textbf{Route}: compute \(\boldsymbol{\alpha}^{(t)}\).
    \item \textbf{Distill}: update student with \(\mathcal{L}^{(t)}\).
    \item \textbf{Gate}: if \(\text{AUC}_m^{(t)} < \theta_{\text{low}}\), decay \(\alpha_m\); if \(\text{AUC}_m^{(t)} > \theta_{\text{high}}\), trigger post-hoc localization.
    \item \textbf{Fallback}: if overlap \(< 0.3\), revert to RF-Domain only.
\end{enumerate}

\section{Experimental Protocol on ChronoBerg}
\label{sec:exp}

\paragraph{Data.} ChronoBerg provides temporally ordered semantic evolution. We partition into \(T\) time windows, treating window \(0\) as reference.

\paragraph{Modalities.} Text embeddings (EmbeddingGemma), audio embeddings (CLAP audio encoder), image embeddings (CLIP ViT).

\paragraph{Baselines.}
\begin{itemize}
    \item \textbf{B1: Static Distillation.} \(\alpha_m = 1/|\mathcal{M}|\) for all \(t\).
    \item \textbf{B2: Covariate-Only Routing.} \(\alpha_m \propto \text{AUC}_m^{(t)}\).
    \item \textbf{B3: AGOD (Ours).} \(\alpha_m \propto \text{AUC}_m^{(t)} \cdot \text{VIMP}_m^{(t)} + \gamma \cdot \text{PO-risk}_m^{(t)}\).
\end{itemize}

\paragraph{Metrics.}
\begin{itemize}
    \item Drift-subgroup Recall@K.
    \item Modality attribution consistency (rank correlation).
    \item Compute cost (FLOPs / wall-clock).
    \item Forgetting rate on stable modalities.
\end{itemize}

\paragraph{Expected Findings.}
AGOD should outperform B1 and B2 on drift-subgroup Recall@K, with the largest gains when concept drift dominates covariate shift. Compute savings should be substantial when only one modality drifts.

\paragraph{CPU prototype.}
The accompanying implementation (\texttt{Python/agod}) ships a ChronoBerg-calendar stream whose shift structure is identifiable: text undergoes covariate shift (semantic evolution), audio undergoes concept drift, and image is stationary. RF-Domain and CFPerm PO-risk are computed with scikit-learn stand-ins for the production CLIP/CLAP/EmbeddingGemma encoders, so the routing law can be unit-tested on CPU.

\section{Discussion: A New Branch of Online Learning}
\label{sec:discussion}

Classical online learning updates parameters from loss gradients. AGOD updates parameters from a \emph{distribution-shift statistic}. This is a distinct branch: the learning signal is not the prediction error alone, but the \emph{structure of the data change}.

\paragraph{Daily Iteration Plan.}
\begin{itemize}
    \item \textbf{Day 1 (Today)}: LaTeX skeleton \& ChronoBerg data loader.
    \item \textbf{Day 2}: Implement RF-Domain \& CFPerm to extract \(g_m^{(t)}\).
    \item \textbf{Day 3}: Implement Softmax routing \& dynamic loss weighting.
    \item \textbf{Day 4}: Baseline comparisons (B1 vs B2 vs B3).
    \item \textbf{Day 5}: Ablation studies (temperature \(\tau\), regularizer \(\lambda\)).
\end{itemize}

\section{Conclusion}
We proposed AGOD, an online distillation framework where modality-specific gap decomposition acts as a clever covariate. By routing distillation capacity toward drifting modalities, AGOD offers a principled alternative to static and covariate-only distillation. We position this as an online learning paradigm where the update signal is a distribution-shift statistic.

\end{document}
